% =====================================================================
%  1bat-ud2-1.tex  —  HORA 1: Del nombre concret a la lletra
%  (sense preàmbul: s'inclou des de main.tex)
% =====================================================================
\horatitol{Sessió 1 — Del nombre concret a la lletra}%
{Traduir situacions quotidianes al llenguatge algebraic: de la frase a la lletra.}

\subsection{Idea clau}

Quan una quantitat encara no la coneixem o pot canviar, la representem amb una
\textbf{lletra} (una \emph{variable}, sovint $x$). Una mateixa frase es pot escriure
amb símbols matemàtics; aquesta traducció és la base de tot el que farem a la unitat.

\begin{idea}
\textbf{Petit diccionari frase $\rightarrow$ símbol}
\begin{itemize}[leftmargin=1.4em, itemsep=2pt]
  \item «el doble de $x$» $\rightarrow$ $2x$ \qquad «la meitat de $x$» $\rightarrow$ $\dfrac{x}{2}$
  \item «$10$ més que $x$» $\rightarrow$ $x+10$ \qquad «$10$ menys que $x$» $\rightarrow$ $x-10$
  \item «el triple de $x$, menys $5$» $\rightarrow$ $3x-5$
  \item «$x$ àpats a $2\eur$ cadascun» $\rightarrow$ $2x$ (cost total)
\end{itemize}
\end{idea}

\subsection{Exemples}

\begin{exemple}[del menjador comunitari]
Un menjador comunitari cobra $2\eur$ per àpat. Si es serveixen $x$ àpats,
la recaptació del dia és
\[
R = 2x \quad(\text{en euros}).
\]
Per a $x=10$: $R=2\per 10 = 20\eur$. \quad Per a $x=120$: $R=2\per 120 = 240\eur$.
La lletra $x$ ens permet escriure \emph{una sola} expressió vàlida per a qualsevol nombre d'àpats.
\end{exemple}

\begin{exemple}[posar nom als elements]
A l'expressió $2x$, el nombre $2$ és el \textbf{coeficient} (què multiplica la lletra)
i $x$ és la \textbf{variable}. A l'expressió $2x+50$, el nombre $50$ que va sol
és el \textbf{terme independent} (no depèn de $x$).
\end{exemple}

\subsection{Exercicis resolts}

\begin{exresolt}[traduir frases a símbols]
Tradueix a llenguatge algebraic, usant $x$ per a la quantitat desconeguda:
\begin{enumerate}[label=\alph*), leftmargin=1.6em]
  \item El doble d'un nombre, augmentat en $7$.
  \item Un monitor cobra $12\eur$ per hora; quant cobra per $x$ hores?
  \item La meitat dels usuaris d'un grup de $x$ persones.
\end{enumerate}
\solucio
\begin{enumerate}[label=\alph*), leftmargin=1.6em]
  \item «el doble d'un nombre» és $2x$; «augmentat en $7$» hi suma $7$:
        \[ 2x+7. \]
  \item Cada hora són $12\eur$, i en fa $x$:
        \[ 12x \quad(\text{en euros}). \]
  \item La meitat de $x$:
        \[ \dfrac{x}{2}. \]
\end{enumerate}
\resposta{(a) $2x+7$ \quad (b) $12x$ \quad (c) $\dfrac{x}{2}$}
\end{exresolt}

\begin{exresolt}[de la situació a l'expressió]
Una entitat reparteix lots d'aliments. Cada lot pesa $3$ kg i, a més, sempre
s'hi afegeix una caixa fixa de $2$ kg de material informatiu per enviament.
Escriu el pes total d'un enviament de $x$ lots.
\solucio
El pes dels lots és $3x$ (kg) i sempre s'hi suma la caixa fixa de $2$ kg:
\[
P = 3x + 2 \quad(\text{en kg}).
\]
Comprovació amb $x=4$ lots: $P = 3\per 4 + 2 = 14$ kg.
\resposta{$P = 3x+2$ kg}
\end{exresolt}

\subsection{Exercicis per fer}

\begin{exercicis}
\begin{exlist}
  \item Tradueix a llenguatge algebraic (usa $x$):
    \begin{enumerate}[label=\alph*), leftmargin=1.6em, topsep=2pt]
      \item El triple d'un nombre.
      \item Un nombre disminuït en $8$.
      \item El doble d'un nombre, menys $5$.
      \item La meitat d'un nombre, més $3$.
    \end{enumerate}
  \item Un casal cobra $4\eur$ per infant i activitat. Escriu el cost total
        si hi participen $x$ infants en una activitat.
  \item Una treballadora familiar fa visites a domicili. Cada visita dura
        $45$ minuts. Escriu, en minuts, el temps total de $x$ visites.
  \item Una associació ven samarretes solidàries a $8\eur$ cadascuna i, a més,
        rep una donació fixa de $100\eur$. Escriu els diners recaptats si ven
        $x$ samarretes.
  \item Escriu \emph{amb les teves paraules} què és una variable i posa un
        exemple propi, de l'àmbit social, d'una frase traduïda a símbols.
  \item \textbf{(Inversa)} Escriu una frase en català que es correspongui amb
        cada expressió: \quad (a) $5x$ \quad (b) $x-12$ \quad (c) $2x+10$.
\end{exlist}
\end{exercicis}

% ---------------------------------------------------------------------
\fullsolucions{Sessió 1}
\begin{solucions}
\begin{sollist}
  \item (a) $3x$ \quad (b) $x-8$ \quad (c) $2x-5$ \quad (d) $\dfrac{x}{2}+3$
  \item $4x$ (euros)
  \item $45x$ (minuts)
  \item $8x+100$ (euros)
  \item Resposta oberta. Exemple vàlid: «el triple d'usuaris atesos» $\rightarrow 3x$.
  \item (a) «el quíntuple d'un nombre» (o «$5$ vegades $x$») \quad
        (b) «un nombre disminuït en $12$» \quad
        (c) «el doble d'un nombre, més $10$»
\end{sollist}
\end{solucions}
